% Source: https://artofproblemsolving.com/community/c3h3267911_2023_lmt_fall_guts_round_p1p15_lexington_mathematical_tournament
% Stable identity: aops-topic-3267911-q06
The condition means $x\in[\dfrac{n}{20},\dfrac{n+1}{20})$, so $23x$ spans a half-open interval of length $\dfrac{23}{20}$. Write its initial fractional part as $\dfrac{j}{20}$, where $j\equiv3n\pmod{20}$ and $0\le j<20$. Three floor values occur exactly when the interval extends strictly past the second integer above its initial floor: $\dfrac{j+23}{20}>2$, or $j\ge18$. The first residues $3n\pmod{20}$ for $n=1,\ldots,6$ are $3,6,9,12,15,18$. Thus the least is $\boxed{n=6}$. Indeed $23x\in[6.9,8.05)$ then allows 6,7,8.
